\SourcePage{347}{350}
\section{Atom in an electric field}

When an atom is placed in an external electric field, its energy levels change; this is the \emph{Stark effect}.

In a uniform external electric field the electrons move in an axially symmetric field, formed by the nucleus and the external field. Strictly, the total angular momentum is no longer conserved; only its projection $M_J$ on the field direction remains conserved. States with different $M_J$ have different energies, so the field removes the directional degeneracy. The removal is incomplete: states differing only in the sign of $M_J$ still have equal energies. Indeed, the atom is symmetric under reflection in any plane through the symmetry axis, chosen below as $z$. States related by such a reflection must have equal energies, while reflection in a plane through an axis reverses the angular momentum about that axis.

Assume the field is weak enough that its additional energy is small compared with separations of neighbouring atomic levels, including fine-structure intervals. The shifts may then be calculated by perturbation theory (\S~38, 39). The perturbation is the energy of the electron system in the uniform field $\boldsymbol{\mathcal E}$:
\begin{equation}\widehat V=-\mathbf d\boldsymbol{\mathcal E}=-\mathcal E d_z,\tag{76.1}\end{equation}
where $\mathbf d$ is the dipole moment. Zeroth-order levels are directionally degenerate, but this degeneracy is immaterial here, and perturbation theory may be applied as for nondegenerate levels. Only matrix elements of $d_z$ for transitions with unchanged $M_J$ can be nonzero (\S~29), so states with different $M_J$ behave independently.

\SourcePage{348}{351}
The first-order shift is given by diagonal elements of the perturbation, but diagonal dipole-moment elements vanish (\S~75). Hence the electric-field splitting is second order in the field.\footnote{Hydrogen is an exception: its Stark effect is linear in the field (see the next section). Other atoms in highly excited, and therefore hydrogen-like, states (\S~68) behave similarly in sufficiently strong fields.} Being quadratic in the field, the shift of level $E_n$ has the form
\begin{equation}\Delta E_n=-\frac12\alpha_{ik}^{(n)}\mathcal E_i\mathcal E_k,\tag{76.2}\end{equation}
where $\alpha_{ik}^{(n)}$ is symmetric. With $z$ along the field,
\begin{equation}\Delta E_n=-\frac12\alpha_{zz}^{(n)}\mathcal E^2.\tag{76.3}\end{equation}

The tensor $\alpha_{ik}^{(n)}$ is also the \emph{polarizability} of the atom. In (11.16), take the parameters $\lambda$ as the components $\mathcal E_i$ and put $\widehat H=\widehat H_0-\mathcal E_i\widehat d_i$. Then the induced mean dipole moment is
\[\overline d_i^{(n)}=-\frac{\partial\Delta E_n}{\partial\mathcal E_i},\]
so that
\begin{equation}\overline d_i^{(n)}=\alpha_{ik}^{(n)}\mathcal E_k.\tag{76.4}\end{equation}
By the general second-order formula (38.10),
\begin{equation}\alpha_{ik}^{(n)}=-2\sum_m'\frac{(d_i)_{nm}(d_k)_{mn}}{E_n-E_m}.\tag{76.5}\end{equation}

The polarizability depends on the unperturbed state, including $M_J$. Its dependence on $M_J$ follows generally by regarding the values $\alpha_{ik}^{(n)}$ as eigenvalues of
\begin{equation}\widehat\alpha_{ik}^{(n)}=\alpha_n\delta_{ik}+\beta_n(\widehat J_i\widehat J_k+\widehat J_k\widehat J_i-\tfrac23\delta_{ik}\widehat{\mathbf J}^{,2}).\tag{76.6}\end{equation}

\SourcePage{349}{352}
This is the general symmetric rank-two tensor depending on $\widehat{\mathbf J}$ (cf. \S~75). Thus
\begin{equation}\Delta E_n=-\frac{\mathcal E^2}{2}\{\alpha_n+2\beta_n[M_J^2-\tfrac13J(J+1)]\}.\tag{76.7}\end{equation}
Summing over all $M_J$ makes the second term vanish; the first therefore shifts the centre of gravity of the split level. For $J=1/2$ the level remains unsplit, consistently with Kramers' theorem (\S~60).

In a nonuniform external field that varies little across the atom, a splitting linear in the field may also arise from the atomic quadrupole moment. The quadrupole interaction operator has the classical form (see II, \S~42)
\begin{equation}\widehat V=\frac16\frac{\partial^2\varphi}{\partial x_i\partial x_k}\widehat Q_{ik},\tag{76.8}\end{equation}
where $\varphi$ is the electric potential and the derivatives are evaluated at the atom.

\ProblemHeading{1.} Determine how the Stark splitting of the components of a multiplet level depends on $J$.

\SolutionHeading Reverse the order in which the perturbations are introduced: first consider Stark splitting without fine structure, then introduce spin--orbit interaction. Since spin does not interact with the external electric field, a level of orbital angular momentum $L$ splits according to (76.2), with
\[\widehat\alpha_{ik}=a\delta_{ik}+b(\widehat L_i\widehat L_k+\widehat L_k\widehat L_i-\tfrac23\delta_{ik}\widehat{\mathbf L}^{,2}).\]
After spin--orbit interaction is introduced, average this operator over states of specified $J$ but unspecified $M_J$, exactly as in problem 1 of \S~75. Equations (76.6), (76.7) result, with
\[\alpha=a,\qquad \beta=b\frac{3(\mathbf{JL})[2(\mathbf{JL})-1]-2J(J+1)L(L+1)}{J(J+1)(2J-1)(2J+3)}.\]
This determines the $J$ dependence, though not the $L,S$ dependence of $a,b$.

\SourcePage{350}{353}
\ProblemHeading{2.} Determine the splitting of a doublet level ($S=1/2$) in an arbitrary, not necessarily weak, electric field.

\SolutionHeading If the splitting is not small relative to the doublet interval, the electric-field perturbation and spin--orbit interaction must be included together:
\[\widehat V=A\widehat{\mathbf S}\widehat{\mathbf L}-\tfrac12\mathcal E^2\{a+2b[\widehat L_z^2-\tfrac13L(L+1)]\}.\]
Dropping terms common to all components,
\[\widehat V=\frac A2[\widehat S_+\widehat L_-+\widehat S_-\widehat L_++2\widehat S_z\widehat L_z]-b\mathcal E^2\widehat L_z^2.\]
For fixed $M=M_J$, form the secular equation in the states $|M_LM_S\rangle=|M\mp1/2,\,\pm1/2\rangle$. Its elements are
\begin{align*}
\langle M-1/2,1/2|V|M-1/2,1/2\rangle&=\frac A2(M-\tfrac12)-b\mathcal E^2(M-\tfrac12)^2,\\
\langle M+1/2,-1/2|V|M+1/2,-1/2\rangle&=-\frac A2(M+\tfrac12)-b\mathcal E^2(M+\tfrac12)^2,\\
\langle M-1/2,1/2|V|M+1/2,-1/2\rangle&=\frac A2[(L+M+\tfrac12)(L-M+\tfrac12)]^{1/2}.
\end{align*}
Hence (problem 1, \S~39)
\begin{equation}\Delta E=-b\mathcal E^2M^2\pm\sqrt{\frac{A^2}{4}(L+\tfrac12)^2+b\mathcal E^2(b\mathcal E^2+A)M^2}.\tag{1}\end{equation}
Terms common to the doublet have been omitted. Both signs apply for $|M|\leqslant L-1/2$. For $|M|=L+1/2$ only one state exists, and with the same additive constant
\begin{equation}\Delta E=(\tfrac A2+b\mathcal E^2)(L+\tfrac12)-b\mathcal E^2(L+\tfrac12)^2.\tag{2}\end{equation}

\ProblemHeading{3.} Determine the quadrupole splitting in an axially symmetric electric field.\footnote{For an arbitrary field see problem 6, \S~103.}

\SolutionHeading Here
\[\varphi_{xx}=\varphi_{yy}=a,\qquad \varphi_{zz}=-2a,\]

\SourcePage{351}{354}
and the other second derivatives vanish. Equation (76.8) becomes
\[\frac a6(\widehat Q_{xx}+\widehat Q_{yy}-2\widehat Q_{zz})=\frac{Qa}{2J(2J-1)}(\widehat{\mathbf J}^{,2}-3\widehat J_z^2),\]
and therefore
\[\Delta E=a\frac{Q}{2J(2J-1)}[J(J+1)-3M_J^2].\]

\ProblemHeading{4.} Calculate the polarizability of hydrogen in its ground state.

\SolutionHeading Spherical symmetry makes $\alpha_{ik}=\alpha\delta_{ik}$, with
\[\alpha=-2e^2\sum_k'\frac{|z_{0k}|^2}{E_0-E_k}.\]
Introduce $\widehat b$ by $z=(m/\hbar)d\widehat b/dt$. Then $z_{0k}=(im/\hbar^2)(E_0-E_k)b_{0k}$ and
\begin{equation}\alpha=\frac{2ime^2}{\hbar^2}\sum_kz_{0k}b_{k0}=\frac{2ime^2}{\hbar^2}(zb)_{00}.\tag{1}\end{equation}
It suffices to know $\widehat b\psi_0$. From (9.2), writing $\widehat b\psi_0=b(\mathbf r)\psi_0$ and using $\widehat H=-\hbar^2\Delta/2m+U(r)$,
\[\tfrac12\psi_0\Delta b+\nabla b\nabla\psi_0=iz\psi_0.\]
With $b=\cos\theta f(r)$ this becomes
\begin{equation}\frac{f''}{2}+\frac{f'}r-\frac f{r^2}+\frac{\psi_0'}{\psi_0}f'=ir.\tag{2}\end{equation}
The solution must keep $f\psi_0$ finite at zero and infinity. For hydrogen's ground state $\psi_0=e^{-r/a_B}/\sqrt\pi$, $a_B=\hbar^2/me^2$, the solution is $f=-ira_B(a_B+r/2)$. Therefore\footnote{The next section obtains this result by another method.}
\[\alpha=\frac{2i}{a_B}(rf\cos^2\theta)_{00}=\frac{2i}{3a_B}(rf)_{00}=\frac92a_B^3.\]

\SourcePage{352}{355}
\ProblemHeading{5.} Calculate the polarizability of an electron in a bound $s$ state in a potential well of range $a$, where $a\varkappa\ll1$, $\varkappa=\sqrt{2m|E_0|}/\hbar$, and $|E_0|$ is the binding energy.

\SolutionHeading The region inside the well may be neglected in $(z\widehat b)_{00}$, and throughout space one may use
\[\psi_0=\sqrt{\frac{\varkappa}{2\pi}}\frac{e^{-\varkappa r}}r.\]
Equation (2) becomes
\[\frac{f''}{2}-\varkappa f'-\frac f{r^2}=ir,\]
whose boundary-condition solution is $f=-ir^2/2\varkappa$. Equation (1) then gives
\[\alpha=\frac{me^2}{4\hbar^2\varkappa^4}.\]
